\section{相対形式復元と楕円ブロックの具体模型}\label{sec:intro}

滑らかなcubicのKuznetsov成分を
\[
 D^b(X)=\langle\Ku(X),\OO_X,\OO_X(1),\ldots,\OO_X(n-2)\rangle
\]
の規約で取る。rotationは
\[
 \rho_X=\mathbf L_{\OO_X}\circ(-\otimes\OO_X(1))
\]
である。すべての圏はdg enhancementと冪等完備化を保って扱う。
可換cubicについては、\((\Ku(X),\rho_X)\) が \(X\) の射影同型類を決定する
\cite[Theorem 1.7]{LZ}。このTorelli定理は可換モデルの一意性を与えるが、
任意の非可換な圏の組がその像に属することを述べない。
幾何側のenhanced同値は、包含・射影と合成してFourier--Mukai核で表されるため、
同定理のFourier--Mukai条件を満たす。
本節以下では、この存在問題のうち、Artin基底での形式復元、
smooth properな相対dg圏のétale局所実現、
および楕円ブロックを用いる具体模型を扱う。
主定理\ref{gdr:thm:generic}の一般点における復元とは、
基底と追加仮定の範囲を区別して適用する。

\subsection{形式的な存在と一意性}

\begin{theorem}[次数付き変形の形式的可換モデル]\label{thm:main}
 \(X=V(F)\subset\PP^{n+1}_{\CC}\) を滑らかなcubic、\(n\ge3\) とする。
 \(R\) を剰余体が \(\CC\) の有限次元可換局所Artin代数とし、
 \(B_R=\bigoplus_{d\ge0}(B_R)_d\) を中心的な \(R\)-代数で、
 \((B_R)_0=R\)、各次数部分が有限自由、かつ指定された次数付き同型
 \[
 B_R\otimes_R\CC\simeq\CC[x_1,\ldots,x_{n+2}]/(F)
 \]
 を持つものとする。次を定義する。
 \[
 \cP_R=D(\QGr B_R)^c,\qquad
 \cK_R=\langle\OO_R,\ldots,\OO_R(n-2)\rangle^\perp\subset\cP_R,\qquad
 \rho_R=\mathbf L_{\OO_R}\circ(1).
 \]
 このとき \((\cK_R,\rho_R)\) はsmooth properな \(R\)-linear dg圏と自己同値の組である。
 あるmarkedな滑らかなcubic変形
 \(X_R\subset\PP^{n+1}_R\) と、markingを保つ同値
 \[
 (\cK_R,\rho_R)\simeq(\Ku(X_R),\rho_{X_R})
 \]
 が存在する。\(X_R\) のmarked射影同型類は一意である。
 これらはArtin底の変更と両立する。
\end{theorem}

ここで \(\QGr\) は全次数付き加群をtorsion加群で割ったGrothendieck圏であり、
上付き \(c\) はcompact対象を表す。Artin底上では
\(\cP_R\) を \(D^b(\qgr B_R)\) 全体に置き換えない。
例えば可換な定数族でも、底 \(R\) が非正則なら両者は異なる。
後述する体上の有限パラメータでは、残余圏を通常のgraded matrix factorization圏、
すなわちOrlov同値で \(D^b(\qgr B)\) に埋め込まれる成分として扱う。
特異なファイバーで、この成分と \(D(\QGr B)^c\) の直交成分を混同しない。

\begin{corollary}[形式族]\label{cor:formal}
 \(B_\hbar\) が次数ごとに \(\CC[[\hbar]]\)-flatな、定理の斉次座標環の形式的変形なら、
 \[
 F_\hbar\in\CC[[\hbar]][x_1,\ldots,x_{n+2}]_3,\qquad F_\hbar\bmod\hbar=F
 \]
 が存在し、各 \(\CC[\hbar]/(\hbar^{r+1})\) 上で残余圏の組を
 \(V(F_\hbar)\) の組と同一視できる。これらの同値は互換的に選べる。
 形式的cubicのmarked同型類は一意である。
\end{corollary}

これは形式的な主張である。係数の収束や、指定された代数的パラメータ空間への
降下、任意の非零パラメータでの可換モデルの存在を含まない。
また \(F_\hbar=F\) を意味しない。

\subsection{代数的族における局所的存在}

定理\ref{thm:etale}は、有限型複素底 \(S\) 上のsmooth properな代数的dg圏の族
\(\cK_S\) と、可逆perfect bimoduleで表された自己同値 \(\rho_S\) を入力とする。
ある点 \(s\) で \((\cK_s,\rho_s)\simeq(\Ku(X),\rho_X)\) なら、
\(s\) のétale近傍 \(T\) 上に滑らかなcubic \(Y/T\) と
\[
 (\cK_T,\rho_T)\simeq(\Ku(Y/T),\rho_{Y/T})
\]
が存在する。従って、可換cubicで実現される複素点はZariski開部分をなす。
これはgraded flatな代数の族から、相対smooth properなdg模型を作る課題を
省略するものではない。

また \(n\ne4\) では \(\HH^2(\Ku(X))=J(F)_3\)、\(\HH^1=0\) であり、
rotationを忘れた形式的復元も成立する（定理\ref{thm:unpolarized-formal}）。
同じ条件下では、代数的dg圏の族についてもrotationを最初から指定せずに
étale局所モデルを作れる（系\ref{cor:etale-unpolarized}）。
これらは局所的な実現の結果であり、rotationなしの大域的な一意性とは分けて扱う。

\subsection{有限パラメータを含む明示的族}

\[
 S_h=\CC\langle x,y,z\rangle/
 ([x,y]-hz^2,\ [y,z]-hx^2,\ [z,x]-hy^2)
\]
と置き、\(T=\sym(xyz)\) を六つの順列の平均、
\(G_h=x^3+y^3+z^3-2h^2T\) とする。
中心変数 \(u_1,\ldots,u_m\) と斉次三次式 \(P(u)\) に対して
\[
 B_{h,P}=S_h[u_1,\ldots,u_m]/(G_h+P(u))
\]
を考える。\(P=\sum_i u_i^3\) の場合を \(B_{h,m}\) と略記する。

\begin{theorem}[一ブロックと任意の中心三次式の可換モデル]\label{thm:global}
 \(m\ge2\)、\(n=m+1\)、\(h\in\CC\)、\(h^6\ne-27\) とする。
 \(P\in\CC[u_1,\ldots,u_m]_3\) のアフィン臨界点が原点だけであると仮定する。
 \(B_{h,P}\) の残余dg圏とrotationは、滑らかなcubic
 \[
 Y_{h,P}=V\!\left(x^3+y^3+z^3+h^2xyz+P(u)\right)
 \subset\PP^{m+2}
 \]
 の \((\Ku(Y_{h,P}),\rho_{Y_{h,P}})\) と同値である。
 この \(Y_{h,P}\) は、同じrotation付き残余圏を持つ滑らかな可換cubicとして
 射影同型を除いて一意である。
\end{theorem}

この定理はsixfoldでは \(m=5\)、sevenfoldでは \(m=6\) を含む。
中心変数どうしの任意の三次混合項を許すが、一つのSklyaninブロックと中心部分が
分離されているという仮定を保持する。\(P=\sum u_i^3\) はその特殊例である。
複数の非可換ブロックについても、対応する可換cubicの残余圏からの完全忠実埋込みを得る
（命題\ref{prop:multiblock-embedding}）。次の部分的な有限位数条件の下では、その像が全体を生成する。

\begin{theorem}[一つの任意位数ブロックを許す多ブロック可換モデル]\label{thm:multiblock-model}
 \(r\ge2\)、\(m\ge0\)、\(n=3r+m-2\) とする。
 各 \(j\) について \(h_j^6\ne-27\) とし、楕円曲線
 \[
 E_{h_j}=V(x_j^3+y_j^3+z_j^3+h_j^2x_jy_jz_j),\qquad
 o_j=[1:-1:0],\quad p_j=[1:-1:-h_j]
 \]
 上の点 \(p_j\) は、高々一つを除いて有限位数であると仮定する。
 すなわち因子を並べ替えた後、\(j\ge2\) にのみ有限位数を仮定し、\(p_1\) の位数には条件を課さない。
 \(m>0\) の場合、\(P\in\CC[u_1,\ldots,u_m]_3\) は原点だけに臨界点を持つとし、
 \(m=0\) の場合は中心部分を省く。
 異なるSklyaninブロックの元を可換とした代数
 \[
 B_{\mathbf h,P}=
 (S_{h_1}\otimes\cdots\otimes S_{h_r})[u_1,\ldots,u_m]
 \big/\left(\sum_jG_{h_j}+P\right)
 \]
 のrotation付き残余dg圏は、滑らかなcubic
 \[
 Y_{\mathbf h,P}=V\!\left(
 \sum_{j=1}^r(x_j^3+y_j^3+z_j^3+h_j^2x_jy_jz_j)+P(u)
 \right)\subset\PP^{3r+m-1}
 \]
 の \((\Ku(Y_{\mathbf h,P}),\rho_{Y_{\mathbf h,P}})\) と同値である。
 この可換モデルは射影同型を除いて一意である。
\end{theorem}

ここで「有限パラメータ」は形式変数でなく実際の複素数であることを意味し、
楕円曲線上の点が「有限位数」であることとは異なる。
定理は一つの無限位数ブロックを許す。
具体的に \(h_1=2\) は無限位数（補題\ref{lem:order-infinite}）、
\(h_j^6=5\ (j\ge2)\) は位数五（補題\ref{lem:order-five}）となる。
従って二ブロックのsixfold、三ブロックのsevenfoldにも、真に無限位数の因子を含む例が得られる。
二つ以上の無限位数ブロックを任意に許す主張は、本定理には含まれない。
証明は第\ref{sec:finite-center}節に置く。

\begin{proposition}[例外パラメータのhomological singularity]\label{prop:bad-intro}
 \(m\ge0\) とする。Fermat中心部分の族 \(B_{h,m}\) で \(h^6=-27\) とすると、\(S_h[u]\) はAS-regularで
 \(G_h+\sum u_i^3\) は中心正則元であるが、そのgraded matrix factorization圏は
 properではなく、さらに
 \[
 \operatorname{hdim}(\qgr B_{h,m})=\infty
 \]
 である。ここでhomological dimensionは、非零のExtが存在する次数の上限である。
 従って、滑らかな射影cubicのKuznetsov成分とは同値にならない。
\end{proposition}

これはAS-regularな周囲、三次の中心正則元、期待されるHilbert series、
および \(\rho^3\simeq[2]\) だけでは、大域的な可換モデルの存在を保証しないことを示す。
本稿の可換点からのPoisson量子化族に限定した、良い有限パラメータ全体での存在問題は別に残る。
一般のAS-regular中心三次商については、smooth properかつCalabi--Yauな例にも
定理\ref{thm:zhang-cy-fourfold-counterexample}の反例がある。

\subsection{先行研究と本稿の位置}

Orlov \cite{Orlov} とCassidy--Conner--Kirkman--Moore \cite{CCKM} は、
中心正則元による商、特異圏、matrix factorization、残余成分の比較を与える。
体上のwindowのdg模型にはCăldăraru--Tu \cite[Theorem 5.8]{CT} を用いる。
本稿の形式的な構成では、これらの体上の同値をArtin底上へ無条件に延長せず、
De Deken--Lowen \cite{DDL} のabelian deformationとcompact generationを用いる。
smoothnessとpropernessの持ち上げはnilpotent descent
\cite{Mathew}による。

残余圏のHochschild cohomologyはBallard--Favero--Katzarkov \cite{BFK}、
幾何学的Kodaira--Spencer写像はHuybrechts--Rennemo \cite{HR} と、
三次元の場合にはJacovskis--Lin--Liu--Zhang \cite{JLLZ} を用いる。
四次元の余分な二方向に対するrotationの作用は
Huybrechts \cite[Proposition 3.12]{Huy} から確認する。
actual dg変形をHochschild形式変形と同一視する際のcurvature問題
\cite{KL}は、Lurie \cite{Lurie} とBlanc--Katzarkov--Pandit
\cite{BKP} の比較写像の単射性で処理する。

明示的族の楕円ブロックはEtingof--Ginzburg \cite{EG} と
Sklyaninの点曲線の公式 \cite{Walton} に依拠する。
本稿では次数シフトの共役、中心変数を加える際の生成性、
有限群による通常次数への降下を証明する。
可換な係数環についてのThom--Sebastiani定理を、
非可換な係数環に対して仮定を確認せず引用することはしない。

有限中心代数の外部積生成については、Elagin--Lunts--Schn\"urer
\cite[Theorem 4.18]{ELS} が、perfect field上のnoetherian J-2 schemesに載る
coherentな代数に対し、ungradedな有界導来圏の生成定理を証明している。
本稿では、\(\ZZ^r\) の任意の粗化次数を保ち、さらに一方の係数代数には
有限中心性を仮定しない形を証明する（定理\ref{thm:one-sided-boxes}）。
\cite[Theorem 4.18]{ELS} は双方の因子をcoherentな構造代数とするため、
任意のnoetherian係数代数を許すここでの主張は、その定理の直接の適用ではない。
新たな入力はBeilinsonの対角分解 \cite{Beilinson} を非可換係数付きで用いる
多項式拡大の補題である。中心のNoether正規化と分離冪等元によるretractを組み合わせ、
斉次中心局所化の核を次元について帰納する。
対象の細分次数へのliftや、係数代数の有限global dimensionは仮定しない。
有限位数THCRの有限中心性 \cite[Corollary 2.3]{SmithTate} をこの片側のみに適用し、
一つの任意位数ブロックを含む多ブロックのMF比較を得る。
He--Ma--Ye \cite{HMY} のtensor productに対する一般化Kn\"orrer周期性は、
中心元の次数が二である非可換quadricとClifford変形を扱う。
本稿の三次中心元への直接の適用ではない。
周囲のnoetherian性については、Chan \cite{ChanMulti} の多重Serre同値と
面・帯による議論を任意個の楕円因子へ展開し、有限位数の仮定なしで証明する。
Keeler \cite{KeelerMulti} のTheorem 2.7は任意個の可換な可逆bimoduleの
NC-amplenessを判定するが、同論文のTheorem 3.3とCorollary 3.5のnoetherian性は
二因子についての主張である。本稿では任意個の積に必要な帰納法を記す。
基礎となる対角分解と有限中心のungraded生成定理は既知である。
一方有限中心の次数付き定式化および本適用の先行研究上の新規性は確認継続中である。

étale局所復元では、Toën--Vaquié \cite{TV} のperfect objectのmoduliと
smoothnessの無限小判定 \cite{Stacks} を用いる。
新たに示すのは、同値bimoduleとrotation intertwinerを同時にパラメータ化したスタックが、
形式的復元定理によって底上smoothになることである。
Pirozhkov \cite{Pirozhkov} およびLin--Zhang \cite{LZ} のTorelliは
二つの可換モデルの一意性に関する結果であり、この局所的な存在の証明とは役割が異なる。
Huybrechts \cite[Section 9]{HuyUpdate} の四次元における一般的な非可換変形の問題と、
ここで扱うrotation付き変形問題とは、比較する構造の水準が異なる。

\paragraph{定理の適用範囲.}
以下ではArtin形式基底、smooth properな相対dg圏、および
有限係数の楕円ブロック族について、異なる仮定の下で結論を証明する。
これらの組合せの先行研究上の初出は未確認であり、
基礎となる既知定理の新規性は主張しない。
可換点からの量子化族のすべての有限パラメータで同じ復元が成立するとは結論しない。

\section{次数付き変形から相対残余圏を作る}\label{sec:relative}

\subsection{nilpotentな底変更によるcompactnessの検出}

\begin{lemma}\label{lem:descent}
 \(R\) を上の局所Artin代数、\(k=R/\fm\) とする。
 \(A\) を \(R\)-linear dg代数、\(M\) を \(A\)-加群とし、
 \(A_0=A\otL_R k\)、\(M_0=M\otL_R k\) と書く。
 \(M_0\) が \(A_0\)-compactなら \(M\) は \(A\)-compactである。
 また導来底変更 \(-\otL_R k\) は保守的である。
\end{lemma}
\begin{proof}
 \(\fm\) の冪による有限濾過から
 \(R\in\thick_{D(R)}(k)\) である。
 \(A\)-加群の任意の族 \((N_\alpha)\) に対し、
 \[
 \bigoplus_\alpha\RHom_A(M,P\otL_RN_\alpha)
 \longrightarrow
 \RHom_A\!\left(M,P\otL_R\bigoplus_\alpha N_\alpha\right)
 \]
 が同値となる \(P\in D(R)\) の全体はthickである。
 \(P=k\) では導来随伴と \(M_0\) のcompactnessによって同値となるので、
 \(P=R\) でも同値となる。保守性も、
 \(P\otL_RM=0\) となる \(P\) のthick部分圏を使って同様に従う。
 これは \cite[Propositions 3.28, 3.35; Remark 3.29]{Mathew}
 の特別な場合である。
\end{proof}

\begin{remark}
 この補題は \(k\) が \(R\)-perfectであるとは仮定していない。
 また普通の商 \(A/\fm A\) だけを調べる命題ではなく、
 必ず導来ファイバー \(A\otL_Rk\) を使う。
\end{remark}

\subsection{abelian deformationと明示的な生成子}

\begin{proposition}\label{prop:qgr}
 定理\ref{thm:main}の \(B_R\) について、
 \(\QGr B_R\) は \(\operatorname{Qcoh}(X)\) のflat abelian deformationである。
 \[
 E_R=\bigoplus_{i=0}^{n+1}\OO_R(i)
 \]
 は \(D(\QGr B_R)\) のcompact generatorであり、
 \[
 \Lambda_R=\REnd(E_R),\qquad
 \Lambda_R\otL_Rk\simeq\REnd_X(E_0)
 \tag{\(\ast\)}\label{eq:derived-basechange}
 \]
 がdg代数として成り立つ。
\end{proposition}
\begin{proof}
 まず \(B_R\) は左右noetherianである。
 \(\fm^aB_R/\fm^{a+1}B_R\) は、
 flatnessにより \(B_0\) の有限個の直和である。
 有限濾過と \(B_0\) のnoetherian性を使えばよい。
 次数付き生成元を持ち上げる同じ議論から、
 \(B_R\) は有限生成で、付随する \(\ZZ\)-代数のtailsも通常のtorsion商で記述できる。

 \(X\) の射影正規性により \(B_0\) の \(\ZZ\)-代数は
 \((\OO_X(i))_{i\in\ZZ}\) のdirected Hom-categoryである。
 この列は \(\ZZ\)-generatingであり、\(i\le j\) に対して
 \[
 \Ext_X^{1,2}(\OO_X(i),\OO_X(j))
 =H^{1,2}(X,\OO_X(j-i))=0.
 \]
 hypersurfaceの完全列と射影空間のcohomologyからこの消滅が従う。
 よって \cite[Theorem 4.5]{DDL} を適用できる。
 \(\OO_R(i)\) は明示的なflat liftである。
 実際 \(B_R(i)\) との有限 \(R\)-加群のテンソル積は完全で、
 \(\operatorname{Gr}B_R\to\QGr B_R\) は完全である。
 これらの対象はabelian圏も生成するので、必要なflat resolutionを取れる。

 閉埋め込み \(j:X\hookrightarrow\PP^{n+1}\) を考える。
 \(E_0=j^*(\bigoplus_{i=0}^{n+1}\OO_{\PP^{n+1}}(i))\) はcompactである。
 \(\RHom_X(E_0,M)=0\) なら、随伴とBeilinsonの生成定理
 \cite{Beilinson}から \(Rj_*M=0\)、従って \(M=0\) となる。
 したがって \(E_0\) はcompact generatorである。

 \(E_R\) は有界flat liftなので、
 \cite[Propositions 5.8, 5.9]{DDL} によりcompact generatorとなる。
 同論文Proposition 5.5と導来随伴は
 \[
 k\otL_R\RHom(E_R,E_R)
 \simeq\RHom_X(k\otL_RE_R,k\otL_RE_R)
 \]
 を与える。同論文Theorem 5.13の証明でも同じ底変更が使われている。
 dg enhancement上の導来還元関手からこの比較を作るため、
 比較は合成と単位を保つ。必要ならhomotopically \(R\)-flatなdg模型を選ぶ。
 これで \eqref{eq:derived-basechange} は複体だけでなくdg代数として成立する。
\end{proof}

\begin{corollary}\label{cor:smooth-proper}
 \(\cP_R=D(\QGr B_R)^c\) はsmooth properであり、
 \(\cP_R\otL_Rk\simeq\Perf(X)\) である。
\end{corollary}
\begin{proof}
 導来Morita理論 \cite{Keller} によって
 \(\cP_R\simeq\Perf(\Lambda_R)\) である。
 \(X\) のpropernessから \(\Lambda_0\) は \(k\)-perfectなので、
 補題\ref{lem:descent}で \(\Lambda_R\) は \(R\)-perfectとなる。
 次に
 \[
 A=\Lambda_R^{\mathrm{op}}\otL_R\Lambda_R,\qquad M=\Lambda_R
 \]
 とする。この導来ファイバーは \(\Lambda_0\) のenveloping algebraと対角加群である。
 \(X\) のsmoothnessにより対角加群はその上でperfectであり、
 同じ補題から \(\Lambda_R\) の対角もperfectとなる。
\end{proof}

\subsection{例外列、mutation、rotation}

\begin{proposition}\label{prop:relative-pair}
 \(\OO_R,\ldots,\OO_R(n-2)\) は \(\cP_R\) 内のrelative exceptional collectionである。
 その右直交 \(\cK_R\) はsmooth properで、導来特殊ファイバーは \(\Ku(X)\) である。
 \(\rho_R=\mathbf L_{\OO_R}\circ(1)\) は \(\cK_R\) の自己同値となる。
 これらの構成はArtin底変更と両立する。
\end{proposition}
\begin{proof}
 Hom複体の導来底変更と補題\ref{lem:descent}から
 \[
 \RHom(\OO_R(i),\OO_R(i))\simeq R,\qquad
 \RHom(\OO_R(j),\OO_R(i))=0\quad(j>i,\ 0\le i,j\le n-2)
 \]
 である。自己Homについては単位写像のconeを還元する。
 propernessにより評価・余評価はperfect対象を保つため、各例外成分はadmissibleである。
 評価のconeによるmutationを順に取って
 \[
 p_R=\mathbf L_{\OO_R}\cdots\mathbf L_{\OO_R(n-2)}
 \]
 と置くと、\(G_R=p_R(E_R)\) は \(\cK_R\) を生成する。
 mutationの導来底変更は \cite[Lemma 5.14]{DDL} による。
 \(G_R\) の導来特殊ファイバーは \(\Ku(X)\) の生成子なので、
 その自己準同型dg代数へ補題\ref{lem:descent}を適用すればsmooth proper性も従う。

 \(T=(1)\) は次数付き射影圏の自己同値で、\(T\OO_R(i)=\OO_R(i+1)\) である。
 \(M\in\cK_R\)、\(1\le i\le n-2\) に対して
 \[
 \RHom(\OO_R(i),TM)=\RHom(\OO_R(i-1),M)=0.
 \]
 \(\OO_R\) によるmutationはこの消滅を保ち、\(i=0\) の直交性も与える。
 よって \(\rho_R\) は \(\cK_R\) を保つ。
 その特殊ファイバーは通常の \(\rho_X\) である。
 Homの比較写像は還元後に同値なので、nilpotentな保守性から
 \(\rho_R\) は完全忠実である。

 \(\rho_R(G_R)\) の特殊ファイバーは生成子である。
 これに直交するInd対象を導来還元すると零になり、
 同じ保守性で元の対象も零になる。従って \(\rho_R(G_R)\) もcompact generatorである。
 冪等完備な圏からの完全忠実完全関手の像はthickなので、\(\rho_R\) は本質的全射でもある。

 底変更については、次数付き導来テンソル関手による生成子の自己Homの自然な比較を使う。
 剰余体上でこの比較は恒等写像となるため、変更後のArtin底上で同値となる。
 評価写像とそのconeも底変更と両立するので、残余圏とrotationまで両立する。
\end{proof}

\section{rotation付きの形式的復元}\label{sec:formal}

\subsection{正確な変形問題}

\begin{definition}
 \((\Ku(X),\rho_X)\) のmarked pair deformationとは、
 smooth properな \(R\)-linear dg圏 \(\cK_R\)、invertible bimoduleで表される
 自己同値 \(\rho_R\)、および特殊ファイバーでのenhanced pairの指定された同定からなる。
 同値にはこのmarkingとの両立を要求する。
 \(\rho_R^3\simeq[2]\) の指定されたcoherenceは定義に含めない。
\end{definition}

\begin{theorem}[pairの形式的復元]\label{thm:pair}
 \(X\) を滑らかなcubic \(n\)-fold、\(n\ge3\) とする。
 marked cubic deformationからmarked pair deformationへの写像は、
 各古典的局所Artin \(\CC\)-代数上で同型類の全単射を与える。
 互換的な形式族についても同型類の一意性と互換的な実現が成り立つ。
\end{theorem}

この主張はfull higher deformation groupoidの同値ではない。
pairには自然変換のscalar automorphismがあり、実現する圏同値も一意とは限らない。

\subsection{actual変形とHochschild形式変形}

\(\cC=\Ind(\Ku(X))\) とする。
actual compactly generated deformationの空間と、その形式moduli problemの間には
\[
 \theta_{\cC}:\operatorname{CatDef}_{\cC}^{c}
 \longrightarrow\operatorname{CatDef}_{\cC}^{\wedge}
\]
がある。右辺の接複体は \(CC^\bullet(\Ku(X))[2]\) である。
actual Morita変形が任意のHochschild classを実現するとは一般にはいえない
\cite{KL}。本稿ではその全射性を使わない。

\begin{lemma}\label{lem:theta}
 上の \(\theta_{\cC}\) は \((-1)\)-truncatedである。
 従って、右辺で同値となった二つのactual変形の同値はactual側へ戻せる。
 この性質はrotationによる \(\ZZ\)-homotopy fixed pointを取った後にも成り立つ。
\end{lemma}
\begin{proof}
 \(\Ku(X)\) はproperでcompact generatorを持つためtameである。
 \cite[Proposition 5.3.21, Theorem 5.1.9]{Lurie} を適用する。
 \cite[Definition 3.3, Section 3.3]{BKP} にも、
 \(\pi_0\) で単射、\(\pi_i\), \(i\ge1\), で同型という形でこの帰結が明記されている。
 比較写像は圏同値に自然であり、markingを移すrotationの作用と両立する。
 空間のmonomorphismはhomotopy limitで保たれる。
\end{proof}

一つのenhanced autoequivalenceを持ち上げる問題は、
markingに作用する \(\ZZ\) のhomotopy fixed pointである。
\(\ZZ\) は一つの可逆射で自由生成されるので、追加の有限群関係を課していない。
\[
 P=\left(\operatorname{CatDef}_{\cC}^{\wedge}\right)^{h\ZZ}
\]
と置くと、これはformal moduli problemであり、
\[
 T_P=R\Gamma(\ZZ,CC^\bullet(\Ku(X))[2]).
\]
\(\ZZ\) の自明加群の長さ一の自由分解から、次の完全列を得る。
\[
 0\longrightarrow
 \operatorname{coker}(1-\rho:\HH^1\to\HH^1)
 \longrightarrow H^0(T_P)
 \longrightarrow
 \ker(1-\rho:\HH^2\to\HH^2)
 \longrightarrow0.                                      \tag{1}\label{eq:tangent}
\]

\subsection{Hochschild計算とKodaira--Spencer写像}

\(N=n+2\) および
\[
 J(F)=\CC[x_1,\ldots,x_N]/(\partial_1F,\ldots,\partial_NF)
\]
を用いる。

\begin{lemma}\label{lem:HH}
 \(n\ge3\) のとき
 \[
 \HH^{<0}(\Ku(X))=0,\quad \HH^0(\Ku(X))=\CC,\quad
 \HH^1(\Ku(X))=0,
 \]
 \[
 \HH^2(\Ku(X))=
 \begin{cases}
 J(F)_3& n\ne4,\\
 J(F)_3\oplus\CC^2& n=4.
 \end{cases}
 \]
 幾何学的Kodaira--Spencer写像は同型
 \[
 H^1(X,T_X)\xrightarrow{\sim}
 \HH^2(\Ku(X))^{\rho_X}=J(F)_3
 \]
 を与える。その次元は \(\binom{n+2}{3}\) である。
\end{lemma}
\begin{proof}
 \cite[Theorem 5.39]{BFK} のKoszul cohomologyによる公式を
 三次のequivariant matrix factorizationへ適用する。
 恒等sectorは偶数次数 \(2r\) に \(J(F)_{3r}\) を与える。
 非自明な \(g\in\mu_3\) では固定点集合は原点だけであり、
 conormalのrankは \(N\)、determinant characterは \(-N\) である。
 cohomological degree \(q\) で残るweightは
 \[
 N+3(q-N)/2=(3q-N)/2.
 \]
 従って二つの非自明sectorは、\(3\mid N\) のときに限り、
 \(q=N/3\) に各一次元を与える。これで列記した群が得られる。

 BFKの恒等sectorのKoszul模型では、\(G\) の乗法が \([G]\) を表す。
 \(F+\varepsilon G\) の実際の変形類との一致は次の比較による。
 \(n\ge4\) では、periodicityから作るkernel morphismと
 Atiyah--Kodaira--Spencer morphismの一致を
 \cite[Lemma 2.19]{HR} とその前後の構成で確認できる。
 ここで三次に必要な次数条件 \(3\le(n+2)/2\) は満たされる。
 \(n=3\) では \cite[Theorem 4.4(2)]{JLLZ} が必要な接写像の同型を与える。
 normal sequenceとEuler sequenceは \(H^1(T_X)=J(F)_3\) を与える。

 幾何学的変形はrotationを持つので、この部分空間は不変である。
 \(n\ne4\) ではそれが \(\HH^2\) 全体である。
 \(n=4\) では \cite[Proposition 3.12]{Huy} により、
 rotationはMukai格子の \(A_2^\perp\) を固定し、\(A_2\) の三つの根を巡回置換する。
 invariantな2-Calabi--Yau volumeによる
 \(\HH^2\simeq\HH_0\) の下で、余分な二方向の固有値は
 \(\zeta,\zeta^2\) となり、不変部分はない。
 最後に偏微分は二次の正則列であるから
 \(H_{J(F)}(t)=(1+t)^N\) であり、次数三の次元は \(\binom N3\) である。
\end{proof}

\subsection{持ち上げの差と自己同型}

\begin{lemma}[Artin帰納法]\label{lem:induction}
 \(D\) をsmoothかつinfinitesimally homogeneousなmarked変形関手で、
 古典的Artin底上ではset-valuedとし、
 \(P\) をformal moduli problemとする。
 自然な写像 \(f:D\to P\) が接空間で同型を誘導すると仮定する。
 さらに、各small extension \(R'\to R\) と \(x'\in D(R')\) に対し、
 \(f(x'\vert_R)\) の任意の自己同型類が \(f(x')\) 上へ持ち上がると仮定する。
 このとき \(D(R)\to\pi_0P(R)\) は各古典的Artin底上で全単射である。
\end{lemma}
\begin{proof}
 Artin長について帰納法を行う。核 \(I\) のsmall extension上に
 \(y'\in P(R')\) があり、その還元を \(f(x)\) と同一視できたとする。
 \(D\) のsmoothnessにより \(x\) を \(x'\) に持ち上げる。
 固定した還元の同定を持つ \(f(x')\) と \(y'\) は、
 \(P(R')\to P(R)\) の同じhomotopy fibreに属する。
 このfibreが空でなければ、その連結成分は
 \(H^0(T_P)\otimes I\) のtorsorである。
 接写像の全射性によって \(x'\) を修正すれば差を消せる。

 単射性では、\(f(x')\simeq f(z')\) の還元を考える。
 帰納法により幾何学的還元は一致するが、
 圏同値の還元はその像の自己同型となり得る。
 仮定によってこの自己同型を上へ持ち上げ、その逆を合成する。
 すると還元を恒等に固定したfibreで二つのliftを比較できる。
 差の接ベクトルの像は零であり、接写像の単射性で元の差も零となる。
\end{proof}

\begin{lemma}\label{lem:automorphisms}
 cubicから得たpairのmarked deformationにおいて、
 small extensionに沿う上の自己同型持ち上げ条件が成り立つ。
\end{lemma}
\begin{proof}
 smooth properな圏ではHochschild複体はperfectで、その導来底変更が成り立つ。
 補題\ref{lem:HH}と局所Artin環上のminimal perfect complexを使うと
 \[
 \HH^0(\cK_R/R)=R,\qquad \HH^1(\cK_R/R)=0
 \]
 である。実際特殊ファイバーの負次数と次数一が零なので、
 minimal complexにはそれらの次数の項がなく、次数零は単位で生成される。

 恒等関手に還元されるinvertible bimoduleは相対対角と同型である。
 perfect bimoduleのliftの差はsmall extensionの核 \(I\) に対して
 \(\Ext^1(\Delta,\Delta)\otimes I=\HH^1(\Ku(X))\otimes I=0\)
 で測られ、対角自身がliftを与えるためである。
 この同型は各段階で指定した特殊ファイバーのmarkingを保つように選ぶ。
 従ってpairの自己同型の基礎となる圏同値は恒等と同型であり、
 残るのは \(\rho_R\Rightarrow\rho_R\) のscalar unitである。
 markingを固定すればこのscalarは \(1+\fm_R\) に入り、必ず \(R'\) へ持ち上がる。
 scalarな二次自己同型も持ち上がる。
 formal moduli problem側の自己同型をactual bimoduleで扱えることは、
 補題\ref{lem:theta}の \(\pi_1\) の同型による。
\end{proof}

\begin{proof}[定理\ref{thm:pair}の証明]
 cubicのmarked変形関手はsmoothである。方程式の係数を持ち上げればよく、
 \(H^1(\OO_X)=H^2(\OO_X)=0\) によって \(\OO_X(1)\) は一意に延長する。
 infinitesimalな射影自己同型もない。実際、二次の偏微分の正則列には
 一次係数の非自明なsyzygyがなく、Euler scalarを除けば \(H^0(T_X)=0\) である。
 これはmarkedな幾何学的変形をset-valuedとして扱うことを正当化する。

 式\eqref{eq:tangent}と補題\ref{lem:HH}により接写像は同型である。
 補題\ref{lem:automorphisms}と補題\ref{lem:induction}を適用し、
 最後に補題\ref{lem:theta}でactual pairの同値へ戻す。
 帰納法では各段階で指定した還元の同値を保てるので、
 \(\CC[\hbar]/(\hbar^{r+1})\) に沿った実現も互換的に選べる。
\end{proof}

\begin{proof}[定理\ref{thm:main}と系\ref{cor:formal}の証明]
 命題\ref{prop:relative-pair}の組へ定理\ref{thm:pair}を適用する。
 互換的な幾何学的liftの方程式の係数の逆極限が \(F_\hbar\) を与える。
 この最後の操作は形式族を作るものであり、非零複素数で評価する操作ではない。
\end{proof}

\begin{corollary}\label{cor:period}
 上で構成したArtin pairについて、指定を伴わない同型
 \(\rho_R^3\simeq[2]\) が存在する。
\end{corollary}
\begin{proof}
 \(\rho_R^3[-2]\) は恒等に還元されるinvertible bimoduleである。
 補題\ref{lem:automorphisms}で使った \(\HH^1\)-rigidityを適用する。
 または可換モデルからこの同型を移してもよい。
 特定のmatrix factorization表示が与えるperiodicity同型とのcoherenceの一致は、
 この系の主張に含めない。
\end{proof}

\begin{remark}
 sevenfoldでは \(\HH^3(\Ku(X))\simeq\CC^2\) があるが、
 上の帰納法は \(\HH^3=0\) を要求しない。
 比較している二つのliftは既に存在し、その差を接空間で修正する。
 higher operationsを無視したformalityの仮定も使っていない。
\end{remark}

\subsection{四次元以外ではrotationを忘れた形式復元}\label{sec:unpolarized-formal}

四次元以外では、形式近傍に限れば、復元の入力からrotationを省くことができる。
ここでいうmarked dg変形は、特殊ファイバーを \(\Ku(X)\) と同定した
smooth properな \(R\)-linear dg圏であり、同値にもこのmarkingとの両立を要求する。

\begin{theorem}[rotationを忘れた形式復元]\label{thm:unpolarized-formal}
 \(X\subset\PP^{n+1}_{\CC}\) を滑らかなcubicとし、\(n\ge3\)、\(n\ne4\) とする。
 任意の古典的局所Artin \(\CC\)-代数 \(R\) について、
 marked cubic変形から \(\Ku(X)\) のmarked smooth proper dg変形への写像は、
 同型類の全単射を与える。従って、任意のそのような \(\cK_R\) は
 一意なmarked cubic変形 \(X_R\) の \(\Ku(X_R)\) と同値である。

 また、固定した \(\cK_R\) 上で、\(\rho_X\) を持ち上げるinvertible bimodule
 \(\rho_R\) が存在し、そのmarked同型類は一意である。
 従ってrotationを忘れる写像
 \[
  \{\text{marked smooth proper pair変形}\}/\simeq
  \ \longrightarrow\
  \{\text{marked smooth proper dg変形}\}/\simeq
 \]
 も全単射である。これらの主張はArtin底変更と互換的な形式族についても成り立つ。
\end{theorem}
\begin{proof}
 \(Q=\operatorname{CatDef}_{\Ind\Ku(X)}^{\wedge}\) と置く。
 補題\ref{lem:HH}の幾何学的Kodaira--Spencer写像は、\(n\ne4\) では
 \[
  H^1(X,T_X)\xrightarrow{\sim}
  \HH^2(\Ku(X))=H^0(T_Q)
 \]
 である。cubicのmarked変形関手はsmoothかつset-valuedである。
 補題\ref{lem:induction}を適用するため、幾何学的な変形の像の自己同型を調べる。

 \(\Ku(X_R)\) の恒等に還元される自己同値を表すinvertible bimoduleは、
 \(\HH^1(\Ku(X))=0\) により相対対角bimoduleと同型である。
 これは補題\ref{lem:automorphisms}と同じsmall extensionの帰納法で従う。
 特殊ファイバー上の同型を指定しても、そこでのscalarは \(R^\times\) へ持ち上がるので、
 同型をmarkingと両立するように選べる。
 従ってmarkedな自己同値の同型類は恒等だけであり、
 必要な自己同型持ち上げ条件は満たされる。
 \(Q\) 側の自己同型についても、補題\ref{lem:theta}の比較写像が
 \(\pi_1\) 上の同型を誘導するため同じことが成り立つ。

 補題\ref{lem:induction}から、cubic変形と \(\pi_0Q(R)\) の全単射が得られる。
 幾何学的変形の像はactual smooth proper変形であり、
 補題\ref{lem:theta}によりactual変形間の同値を検出できる。
 よって最初の主張が従う。この議論はHochschild形式変形の一般的なactual実現性を
 仮定せず、この場合には幾何学的なliftで全てを実現している。

 得られた同値 \(\Ku(X_R)\simeq\cK_R\) を通じて
 \(\rho_{X_R}\) を移せば、\(\rho_X\) のliftが存在する。
 二つのmarkedなlift \(\rho_R,\rho'_R\) があれば、
 \(\rho'_R\rho_R^{-1}\) は恒等に還元されるinvertible bimoduleである。
 再び \(\HH^1=0\) による対角のrigidityを用いると、
 \(\rho'_R\rho_R^{-1}\simeq\mathrm{id}\) となり、markingを保つ同型も選べる。
 これがliftの同型類の一意性とforgetful mapの単射性を与える。
 各small extensionで還元上の同値を保って構成できるため、
 形式族での互換性も従う。
\end{proof}

\begin{remark}\label{rem:unpolarized-scope}
 この定理の一意性は同型類についてのものである。
 rotationとのintertwinerにはscalar自由度があり、
 rotationを忘れる写像がfull higher deformation groupoidの同値だとは主張しない。
 実際、\(\HH^0=\CC\)、\(\HH^1=0\) とrotationのscalar上の自明作用から、
 pairの接複体には \(H^{-1}(T_P)\simeq\CC\) がある一方、
 \(H^{-1}(T_Q)=0\) である。

 \(n=4\) では幾何学的な20方向に加えて \(\HH^2\) に二方向があり、
 rotationの固定部分を取って初めて接写像が同型となる。
 この次元でrotationを省く結論を上の証明から導くことはできない。
 ここでは、Hochschildの全22方向が各Artin底上でactual変形として
 実現されるとは主張していない。

 \(n=7\) では \(\HH^3\simeq\CC^2\) が残るが、
 幾何学的なliftと既に存在するliftを比較する帰納法により結論は成り立つ。
 障害を収め得るcohomology群が非零であることと、
 この変形問題で実際に障害が現れることは同じではない。
 一般の有限パラメータの圏からrotationを省けるという大域的主張も含めない。
\end{remark}

